\section{方法}

\paragraph{数据收集。}
研究共采用两个相互独立的数据集对模型进行训练与验证，以充分评估模型架构的可靠性，实现针对纳米硒的细胞毒性预测。

\begin{enumerate}
  \item \textbf{通用数据集}：
该数据集来源于 Masarkar 等人 \cite{masarkar2024} 开展的纳米颗粒细胞毒性元分析研究。
数据集汇总了已发表文献中关于金属、金属氧化物、聚合物以及碳基纳米颗粒四大类别的细胞毒性实验数据，是当前纳米材料机器学习毒性预测领域具有代表性的公开数据资源之一。
由于本研究的模型改进工作是在其原有框架基础上展开，因此采用该数据集对基线模型进行复现，以保证模型优化过程具备良好的可比性与有效性。
该数据集原始共包含 4111 条实验记录，经清洗后保留 4102 个有效样本用于后续训练测试划分与模型比较。

  \item \textbf{自建数据集}：
针对纳米硒这一特定研究对象，本研究从 43 篇发表于 2008--2022 年的同行评议文献中系统提取实验记录，构建纳米硒细胞毒性预测专项数据集。数据整理采用大模型粗提取与人工精校验证相结合的流程：首先利用提示词工程，引导大模型按照 Masarkar 等人\cite{masarkar2024} 的特征体系提取 43 篇文献中的特征属性；对于正文未明确报告而可从图表获得的数据，进行人工数字化提取；最后由研究人员逐项对照原始文献核验全部记录。来源文献完整清单见补充材料表~\ref{tab:source-publications}，逐文献记录数尚待根据原始 383 条数据核实。整理后共包含 383 条有效记录，每条记录对应 37 个特征参数，覆盖纳米硒理化性质、暴露条件及生物学终点，为后续增强预测模型的训练、验证与评价提供数据基础。
\end{enumerate}

为进一步说明自建数据集的文献来源结构，本研究按题名和研究目的将 43 篇纳米硒相关文献划分为四类：材料构建/表面修饰、机制毒理/安全性、递送/联合治疗、诊疗/免疫及其他应用。该分类仅用于展示文献覆盖范围和主题分布，不改变后续建模中对原始实验记录的提取规则。图~\ref{fig:senp-literature-cloud} 展示文献发表年份与主题分类的点云分布。

\begin{figure}[H]
  \centering
  \begin{tikzpicture}
    \begin{axis}[
      acs axis,
      width=0.82\linewidth,
      height=6.5cm,
      xmin=2007.5,
      xmax=2022.8,
      ymin=0.55,
      ymax=4.45,
      xlabel={发表年份},
      xtick={2008,2010,2012,2014,2016,2018,2020,2022},
      ytick={1,2,3,4},
      yticklabels={材料/表面,机制/安全,递送/联合,诊疗/免疫/其他},
      y tick label style={font=\scriptsize},
      title={\textbf{(A)} 43 篇文献点云}
    ]
      \addplot+[only marks, mark=*, mark size=1.9pt, acsblue] coordinates {
        (2008,0.78) (2011,1.10) (2011.88,0.92) (2012.12,1.24)
        (2012.36,0.80) (2013.16,1.08) (2014,0.92) (2019,1.20)
      };
      \addplot+[only marks, mark=square*, mark size=1.8pt, acsred] coordinates {
        (2011.15,1.80) (2012.86,2.08) (2013.05,1.86)
        (2013.30,2.22) (2018,1.92) (2019.15,2.12)
      };
      \addplot+[only marks, mark=triangle*, mark size=2.1pt, acsorange] coordinates {
        (2012.44,2.76) (2012.68,3.18) (2013.42,2.90) (2014.04,3.24)
        (2014.32,2.78) (2015,3.12) (2015.96,2.86) (2016.22,3.22)
        (2017,2.94) (2018.16,3.16) (2019.30,2.84) (2020,3.12)
        (2021,2.90) (2021.24,3.20)
      };
      \addplot+[only marks, mark=diamond*, mark size=2pt, acsgray] coordinates {
        (2016.34,3.78) (2016.54,4.06) (2016.74,4.30) (2018.30,3.88)
        (2018.50,4.18) (2019.44,3.78) (2019.64,4.08) (2019.84,4.30)
        (2020.14,3.92) (2021.34,4.20) (2021.54,3.84) (2022,4.10)
        (2022.18,3.78) (2022.36,4.30) (2022.54,3.96)
      };
    \end{axis}
  \end{tikzpicture}

  \vspace{0.8em}

  \begin{tikzpicture}
    \begin{axis}[
      acs xbar,
      width=0.62\linewidth,
      height=4.6cm,
      xmax=16,
      xlabel={文献数量},
      symbolic y coords={机制/安全,材料/表面,递送/联合,诊疗/免疫/其他},
      ytick=data,
      y tick label style={font=\scriptsize},
      title={\textbf{(B)} 分类数量}
    ]
      \addplot+[fill=acsblue, draw=none] coordinates {(6,机制/安全) (8,材料/表面) (14,递送/联合) (15,诊疗/免疫/其他)};
    \end{axis}
  \end{tikzpicture}
  \caption{自建纳米硒数据集 43 篇来源文献的发表年份点云与主题分类分布。}
  \label{fig:senp-literature-cloud}
\end{figure}

\paragraph{数据预处理。}
对于通用数据集，我们沿用了 Masarkar 等人 \cite{masarkar2024} 的预处理方案：
首先，我们删除了数据集中包含大量缺失值的特征属性（如多分散指数 PDI 和长径比），以及与预测无关的文献信息（如文章 ID、发表年份和 DOI 编号），将特征数量从初始的 30 个精简至 15 个。
其次，对于某些在物理化学性质（如尺寸、Zeta 电位、表面电荷和 PDI）上存在严重信息缺失的纳米颗粒（例如量子点 QD 和羟基磷灰石 HAP），我们剔除了其对应的所有数据点。

随后，我们将剩余的数据集按照材料类型统一划分为四大类纳米颗粒：金属（M）、金属氧化物（MO）、聚合物（P）和碳基（C）纳米颗粒。
在处理具体的物理化学特征时，为了整合不同测试条件下的数据，我们采用了以下优先级策略：

\begin{enumerate}
  \item \textbf{尺寸（Size）特征}：
优先采用无任何介质条件下的尺寸；若该值缺失但提供了水中和培养基中的尺寸，则采用水中的尺寸；若原数据仅提供了一种条件（无介质、水中或培养基中）下的尺寸，则直接使用该已知值。

  \item \textbf{Zeta 电位特征}：
采用类似策略，优先采用水中的 Zeta 电位；若仅提供了一种条件下的 Zeta 电位，则直接使用该值。
\end{enumerate}

在执行上述整合、审查目标变量缺失并严格过滤异常值后，最终清洗后的数据集保留了 4102 个有效数据点和 15 个特征属性。
最后，为了使数据完全适合模型训练，我们使用“min-max”缩放器（MinMaxScaler）对编译后的数据集进行了归一化处理，并对所有的类别属性进行了标签编码（Label encoding）。

对于自建的纳米硒专项数据集，预处理步骤进行了改进。
本研究自建的纳米硒毒性数据集共包含 383 个样本与 37 个特征。
数据集覆盖了纳米硒研究中常见的实验条件区间，包括不同的暴露剂量、粒径范围、形态类型与表面修饰策略，可反映纳米硒与其理化性质的关系在多个维度上的变化。

自建数据集数据预处理包括以下步骤：

\begin{enumerate}
  \item \textbf{缺失值处理}：由于粒径、Zeta 电位、暴露时间、浓度、细胞存活率诸字段都有若干缺失值，故首先直接剔除含缺失值的记录。此外，“检测方法”字段中若干条目“Characterization”“Western”“Cellular Uptake”属于非毒性检测目的条目，故一并删除。最终共剔除 262 条记录，保留 383 个符合要求的有效样本。
  \textbf{上述剔除仅针对所列关键字段；派生特征（如粒径比）和非关键分类变量（如表面电荷）的缺失值在剔除之后单独处理，不在此列。}

  \item \textbf{多报告值的优先级选择}：参照 Masarkar 等人 \cite{masarkar2024} 提出的方法：同一纳米颗粒在文献中报告多种尺寸时，按无介质中的粒径 > 水相粒径 > 培养基中的粒径的顺序选取，而 Zeta 电位宜取水相测量值，若无该数据则取文献中所给的唯一测量值。

  \item \textbf{单位统一与换算}：把各浓度单位都统一换算为 $\mu$g/mL，凡以 $\mu$M 为单位的浓度都根据硒的分子量（78.96 g/mol）换算为 $\mu$g/mL，同时把暴露时间统一转换为小时（h）。

  \item \textbf{数值特征变换}：由于浓度、时间两变量都呈偏态分布，对二者做对数变换：
        \[
          \text{Log\_Conc} = \ln(1 + \text{Concentration\_μg/mL}), \quad
          \text{Log\_Time} = \ln(1 + \text{Time\_h})
        \]
        同时引入“粒径比”= 水合粒径 / 核心粒径来表征颗粒在水相中的聚集倾向，所有缺失值均按 1.5 填充。

  \item \textbf{分类变量编码}：对表面电荷（正/负/中性）、细胞类型、修饰基团、检测方法、晶型（结晶/无定形）等类别特征，采用标签编码转换为数值，其中表面电荷缺失值填充为“未知”。

  \item \textbf{相关性检查}：对数值特征计算 Pearson 相关系数，未发现 $r > 0.8$ 的高相关度特征对，表明特征冗余度很低。
\end{enumerate}

\paragraph{特征工程与多模态融合。}

在此基础上，模型输入特征被组织为基础特征和增强特征两组。基础特征包含以下几类信息：暴露条件（浓度、时间）、颗粒理化属性（粒径、Zeta 电位、尺寸比、表面电荷）、生物体系属性（细胞类型）以及实验条件（检测方法、修饰基团）。
针对硒纳米颗粒专用建模，在基础特征之外进一步引入 5 项增强变量，即 IC50、元素 Se 含量、长期稳定性、Se 3d 结合能 和晶型编码。其中，元素 Se 含量、长期稳定性、Se 3d 结合能 与晶型编码主要用于表征材料本征化学属性；IC50 则与毒理响应终点具有较强关联，因此本文将其视为辅助变量，用于比较含/不含该变量时模型表现的差异，而不将其简单等同于可在真实预测前稳定获得的先验理化特征。

为利用不同特征模态之间的互补信息，本研究采用多模态特征融合策略，将输入划分为三个模态：（1）基础理化描述符，包括浓度、时间、粒径、Zeta 电位、尺寸比、表面电荷、细胞类型、检测方法和修饰基团；（2）材料化学增强描述符，包括元素 Se 含量、长期稳定性、Se 3d 结合能和晶型编码；（3）可选辅助描述符 IC50。融合分为三个阶段：首先拼接基础特征与材料化学增强特征，构建联合表示；随后，在 IC50 可用时以其作为查询，对联合表示计算交叉注意力权重，自适应突出相关特征，而在 IC50 缺失时旁路注意力机制、直接使用联合表示；最后加入残差连接以保留原始联合表示，并将所得多模态特征矩阵输入预测模型。该设计将材料本征化学信息与通用理化描述符整合，并根据辅助信息是否可用动态调整特征权重。

为评估新增增强信息对模型性能的提升作用，保持基础特征不变，对上述 5 项增强变量的全部非空组合进行穷举搜索，共 31 种组合，并以三折交叉验证 MAE 作为筛选指标，用于确定表现较优的增强变量子集。需要指出的是，若最优组合包含 IC50，则其结果更适合作为辅助分析依据，而不宜直接等同于真实部署情境下的预测能力提升。

\paragraph{模型与评价。}
本研究围绕表格型纳米毒性数据比较了树模型、注意力增强树模型与深度表格模型三类方法。在基准数据集上，纳入随机森林、XGBoost、ABRF、LARF、TabTransformer 和蒸馏 TabTransformer 六种模型，以系统评估不同架构对异质纳米颗粒数据的适应性；在硒纳米颗粒专用数据集上，重点比较基础特征与增强特征条件下各模型的性能差异。随机森林作为核心基线模型，用于评估中小样本场景下的稳健性与可解释性；XGBoost 用于测试梯度提升树在增强特征条件下的高精度建模能力；ABRF 与 LARF 用于考察注意力机制引入后是否能够提升树模型表现；TabTransformer 与蒸馏 TabTransformer 则作为深度表格模型代表，用于与传统树模型进行对照。

本文对模型性能采用回归、分类两类指标进行评价。回归任务中采用决定系数（\(R^2\)）、平均绝对误差（MAE）和均方根误差（RMSE）。\(R^2\) 是对模型所拟合的目标变量方差的解释比例的自然度量，理论上不超过 1（当模型表现差于均值基线时可能为负值），越接近 1 表明拟合效果越好。MAE 是预测值与真实值绝对误差的算术平均，其单位与目标变量一致（\%），因此也是预测典型偏差的直接度量。RMSE 对较大误差更敏感，适用于考察预测的稳定性。

分类任务沿用 Masarkar 等人\cite{masarkar2024} 所采用的 85\% 细胞存活率阈值，将样本划分为“高存活”和“低存活”两类，以保持与基准数据集一致并便于直接比较模型表现，继而用准确率（Accuracy）、F1 分数和 ROC 曲线下面积（AUC）作为主要指标。准确率即正确分类的比例，F1 分数是精确率、召回率的调和平均，对类别不平衡问题具有适应性，而 AUC 衡量模型区分正负类的能力，0.5 为随机猜测，1.0 表示完全区分。

对于大规模多类型纳米颗粒基准数据集，采用 7:3 训练测试划分，以测试集性能作为主要比较依据。回归任务使用 $R^2$、MAE 和 RMSE 进行评价，分类任务使用 Accuracy、F1 和 AUC 进行评价。

对于硒纳米颗粒专用数据集，采用三折交叉验证，以 CV MAE 和全量 $R^2$ 作为主要回归评价指标，并用于比较基础特征与增强特征两种设置下不同模型的表现。增强特征筛选阶段同样采用三折交叉验证 MAE 作为统一指标，以保证不同特征组合之间的结果可比性。

\paragraph{预测模型的架构细节。}
为系统评估不同架构在纳米毒性预测任务上的适用性，我们尝试了多种模型，涵盖传统树模型、梯度提升树、注意力增强树模型以及深度表格模型。
经过初步性能比较，最终选择随机森林、XGBoost、ABRF、LARF、TabTransformer 以及蒸馏 TabTransformer 六种模型作为主要分析对象。
下面将依次阐述各模型的设计思想及其在本任务中的定位。

随机森林是以决策树为基础的典型的集成学习方法，通过自助采样（Bootstrap）构建多个决策树，并基于其结果进行回归或分类。该方法具有较强的抗过拟合能力，能有效处理混合型特征（数值型和类别型），并且对缺失值不敏感，能直接估计特征重要性，因此在中小型表格数据任务中呗广泛用作基线模型。
基于此，本文采用随机森林来考察其在样本量有限的硒纳米颗粒数据集上的稳健性及可解释性。


XGBoost 是一种基于梯度提升树的集成学习框架，在模型构建过程中引入了二阶导数信息以及 L1、L2 正则化项，因此在表格数据任务中具有较高的预测性能。该方法通过串行迭代方式不断拟合残差，通常能够获得较高的模型精度，但同时对超参数设置较为敏感，需要通过合理调参降低过拟合风险。
本文引入 XGBoost 模型，旨在分析梯度提升树方法在加入元素组成、稳定性、结合能等增强信息后，对硒纳米颗粒数据集预测性能的变化，并进一步比较含/不含 IC50 时模型表现的差异。

ABRF 是在标准随机森林基础上加入了一个注意力网络来对各棵决策树的预测做动态加权处理：先训练普通的随机森林，再对每个样本收集所有树的预测值，接着用一个小型神经网络直接、自动地学习各棵树就该样本的权重，最后用各树预测值的加权和作为最终预测结果。因此其本质是为了用注意力机制合理地挖掘不同树、不同样本之间预测可靠性的差异。
ABRF 是一种在标准随机森林基础上引入注意力机制的集成学习方法。该方法首先训练传统随机森林模型，并收集各决策树对样本的预测结果；随后利用小型神经网络学习不同决策树在不同样本上的权重分配；最终通过对各树预测结果进行加权求和得到模型输出。核心思想是在样本层面建模不同决策树预测贡献的差异，以提升模型的表达能力。
本文引入 ABRF 模型旨在分析单层注意力机制对树模型预测性能的影响。


LARF 在 ABRF 的基础上进一步引入了特征级注意力机制。在模型训练过程中，首先学习特征权重向量，对原始特征进行加权处理，以增强重要特征在决策树分裂过程中的作用；随后结合树级注意力机制，对不同决策树的预测结果进行动态加权，从而形成包含特征级与树级的双层注意力结构。该方法能够同时优化特征表示与树模型集成过程。
本文引入 LARF 来分析更复杂的注意力机制对纳米硒数据集预测性能的影响。

TabTransformer 是一种面向表格数据的深度学习模型，其核心思想是利用 Transformer 的注意力机制对类别特征进行编码，将其表示为嵌入向量，并与数值特征拼接后输入多层感知机进行预测。由于 Transformer 能够建模类别特征之间的高阶交互关系，因此该模型在包含多种类别变量的表格数据任务中具有一定优势。本文将 TabTransformer 作为深度表格模型的代表，用于评估深度学习方法在纳米毒性数据集上的预测性能，并与传统树模型进行比较。

蒸馏 TabTransformer 则基于知识蒸馏方法，以完整的 TabTransformer 作为教师模型，对参数规模较小的学生模型进行训练。该方法通过降低嵌入维度及多层感知机隐藏层规模，实现模型轻量化，并在训练过程中同时利用教师模型输出的软标签与真实标签信息，以兼顾模型精度与计算效率。
故本文引入该模型，旨在分析资源受限场景下轻量化深度表格模型的部署可行性及其性能表现。

由于六种模型在大规模、多类型纳米颗粒基准数据集上已经进行了较为系统的比较，因此本文重点分析不同模型架构对异质纳米颗粒数据的适应能力。
在纳米硒专用数据集上，由于样本量较小（383 条），本文主要选取随机森林、XGBoost 和 TabTransformer 三种代表性模型进行比较，并进一步分析基础特征与增强特征两种特征设置下的性能差异。其中，随机森林与 XGBoost 用于评估传统树模型的表现，TabTransformer 用于考察深度学习模型在小样本表格数据中的适用性。
对于 ABRF、LARF 以及蒸馏 TabTransformer，由于其模型结构相对复杂、训练成本较高，因此仅在大规模基准数据集上进行评估，以兼顾实验复杂度、计算成本与结果分析的可读性。

\paragraph{可视化特征效应与交互。}
特征重要性分析能够确定对模型预测贡献较大的变量，但仅凭重要性排序难以说明特征与细胞存活率之间的关系形态，例如这种关系是近似线性、单调变化还是存在更复杂的非线性响应。参考 Masarkar 等人 \cite{masarkar2024} 的处理思路，本研究先对纳米硒数据集的样本组成、连续型变量和类别型变量进行可视化描述，以明确模型训练前的数据覆盖范围与潜在偏倚；在完成模型训练后，再结合特征重要性进一步考察单个特征效应与双特征交互关系。对于单特征效应，可通过部分依赖图（partial dependence plot, PDP）和个体条件期望图（individual conditional expectation, ICE）观察模型预测值随某一特征变化的整体趋势及样本间异质性；对于两个特征的联合影响，可进一步采用二维 PDP 等高线图展示其组合变化对预测细胞存活率或毒性类别的影响。

在进行特征效应可视化之前，首先需要明确样本组成和变量分布。本研究数据集共包含 383 条有效记录（剔除粒径、Zeta 电位等关键理化性质缺失以及细胞存活率缺失的记录后），来源于 43 篇已发表文献，文献发表年份覆盖 2008--2022 年。清洗后的有效记录主要集中于 2018 年（157 条，占 41.0\%），其次为 2019 年（96 条，占 25.1\%）和 2012 年（44 条，占 11.5\%）。所有纳米硒颗粒的形貌均为球形，物种来源以人类细胞为主（371 条，占 96.9\%）。

图~\ref{fig:senp-year-distribution} 显示了样本来源年份的集中程度，表~\ref{tab:senp_feature_statistics} 汇总了主要连续型特征的描述性统计量，并为后续可视化解释提供了基础。粒径的分布明显右偏（偏度 3.92），绝大多数颗粒集中在 30 nm 附近，极少数样本超过 200 nm，符合纳米硒常规粒径范围。水合粒径右偏程度较轻（偏度 1.69），其均值（105 nm）明显大于核心粒径均值（65 nm），这表明水化层以及可能存在的弱聚集效应。Zeta 电位分布接近对称（偏度 -0.08），范围从 $-60$ mV 到 $+36$ mV，其中负电位占主导，反映出多数表面修饰（如聚糖、PVP）带负电。

\begin{figure}[H]
  \centering
  \begin{minipage}{0.48\linewidth}
    \centering
    \begin{tikzpicture}
      \begin{axis}[
        acs xbar,
        width=\linewidth,
        height=5.1cm,
        xmax=175,
        xlabel={样本量},
        symbolic y coords={2012,其他年份,2019,2018},
        ytick=data,
        y tick label style={font=\scriptsize},
        title={\textbf{(A)} 来源年份频率}
      ]
        \addplot+[fill=acsblue, draw=none] coordinates {(44,2012) (86,其他年份) (96,2019) (157,2018)};
      \end{axis}
    \end{tikzpicture}
  \end{minipage}
  \hfill
  \begin{minipage}{0.48\linewidth}
    \centering
    \begin{tikzpicture}
      \begin{axis}[
        acs axis,
        width=\linewidth,
        height=5.1cm,
        ymin=0,
        ymax=105,
        ylabel={累计占比（\%）},
        symbolic x coords={2018,2019,其他年份,2012},
        xtick=data,
        x tick label style={rotate=35, anchor=east, font=\scriptsize},
        ymajorgrids=true,
        title={\textbf{(B)} 累计贡献}
      ]
        \addplot+[acsred, mark=square*, mark size=1.7pt] coordinates {(2018,41.0) (2019,66.1) (其他年份,88.6) (2012,100.0)};
      \end{axis}
    \end{tikzpicture}
  \end{minipage}
  \caption{纳米硒数据集中文献来源年份的样本量分布与累计贡献。}
  \label{fig:senp-year-distribution}
\end{figure}

浓度分布极端右偏（偏度 8.28），四分位距仅为 $0.52$--$10$ $\mu$g/mL，但最大值高达 $490$ $\mu$g/mL，提示少数实验采用极高剂量，可能对模型产生杠杆效应。暴露时间呈轻微左偏（-0.42），中位数 72 h，说明多数实验采用 24 h 或 72 h 的暴露方案。细胞存活率分布近似均匀（偏度 -0.17），均值与中位数均接近 50\%，意味着数据集正负样本较为平衡，有利于分类模型训练。

为直观展示上述分布，我们进一步基于描述性统计绘制了连续型特征的可用样本量和偏度对照图（图~\ref{fig:senp-continuous-summary}）。该图不对分位数之外的个体观测分布进行外推，而是直接使用表~\ref{tab:senp_feature_statistics} 中的 $n$ 和偏度刻画数据覆盖度与分布偏斜程度。

\begin{table}[H]
  \centering
  \caption{数据集连续型特征的描述性统计（各特征可用样本量见 $n$ 列）}
  \label{tab:senp_feature_statistics}
  \begin{tabular}{lccccccc}
    \hline
    特征 & $n$ & 均值 & 标准差 & 中位数 & Q1 & Q3 & 偏度 \\
    \hline
    粒径 (nm)            & 339 & 65.30  & 79.72  & 30.00  & 30.00  & 92.15   & 3.92  \\
    水合粒径 (nm)        & 279 & 104.97 & 46.73  & 91.00  & 84.00  & 110.00  & 1.69  \\
    Zeta 电位 (mV)       & 144 & $-5.19$ & 29.76  & $-7.50$ & $-29.10$ & 29.10 & $-0.08$ \\
    暴露时间 (h)         & 363 & 55.21  & 26.14  & 72.00  & 24.00  & 72.00   & $-0.42$ \\
    浓度 ($\mu$g/mL)     & 363 & 12.59  & 38.17  & 0.90   & 0.52   & 10.00   & 8.28  \\
    细胞存活率 (\%)      & 371 & 49.43  & 28.50  & 50.00  & 41.00  & 68.00   & $-0.17$ \\
    \hline
  \end{tabular}
\end{table}

\begin{figure}[H]
  \centering
  \begin{minipage}{0.49\linewidth}
    \centering
    \begin{tikzpicture}
      \begin{axis}[
        acs ybar,
        width=\linewidth,
        height=5.8cm,
        ymin=0,
        ymax=405,
        ylabel={可用样本量 $n$},
        symbolic x coords={粒径,水合粒径,Zeta电位,暴露时间,浓度,细胞存活率},
        xtick=data,
        x tick label style={rotate=45, anchor=east, font=\scriptsize},
        title={\textbf{(A)} 数据覆盖度}
      ]
        \addplot+[fill=acsblue, draw=none] coordinates {(粒径,339) (水合粒径,279) (Zeta电位,144) (暴露时间,363) (浓度,363) (细胞存活率,371)};
      \end{axis}
    \end{tikzpicture}
  \end{minipage}
  \hfill
  \begin{minipage}{0.49\linewidth}
    \centering
    \begin{tikzpicture}
      \begin{axis}[
        acs axis,
        width=\linewidth,
        height=5.8cm,
        ymin=-1,
        ymax=9.2,
        ylabel={偏度},
        symbolic x coords={粒径,水合粒径,Zeta电位,暴露时间,浓度,细胞存活率},
        xtick=data,
        x tick label style={rotate=45, anchor=east, font=\scriptsize},
        title={\textbf{(B)} 分布偏斜}
      ]
        \addplot+[only marks, acsred, mark=square*, mark size=1.8pt] coordinates {(粒径,3.92) (水合粒径,1.69) (Zeta电位,-0.08) (暴露时间,-0.42) (浓度,8.28) (细胞存活率,-0.17)};
        \addplot+[black!55, dashed, mark=none] coordinates {(粒径,0) (细胞存活率,0)};
      \end{axis}
    \end{tikzpicture}
  \end{minipage}

  \vspace{0.9em}

  \begin{minipage}{0.62\linewidth}
    \centering
    \begin{tikzpicture}
      \begin{axis}[
        acs axis,
        width=\linewidth,
        height=5.2cm,
        xmin=110,
        xmax=390,
        ymin=-1,
        ymax=9.2,
        xlabel={可用样本量 $n$},
        ylabel={偏度},
        title={\textbf{(C)} 覆盖度--偏度诊断}
      ]
        \addplot+[
          only marks,
          mark=*,
          mark size=1.9pt,
          acsorange,
          nodes near coords,
          nodes near coords style={font=\tiny, anchor=west, xshift=1.5pt},
          point meta=explicit symbolic
        ] coordinates {
          (339,3.92) [粒径]
          (279,1.69) [水合粒径]
          (144,-0.08) [Zeta电位]
          (363,-0.42) [暴露时间]
          (363,8.28) [浓度]
          (371,-0.17) [细胞存活率]
        };
        \addplot+[black!55, dashed, mark=none] coordinates {(110,0) (390,0)};
      \end{axis}
    \end{tikzpicture}
  \end{minipage}
  \caption{主要连续型特征的可用样本量与偏度分布。}
  \label{fig:senp-continuous-summary}
\end{figure}

类别型特征的频率分布同样用于辅助判断模型可能学习到的主要数据模式（图~\ref{fig:senp-categorical-frequency}）。从目前所考察的类别型特征分布来看，数据集涉及 42 种细胞系，结肠癌细胞 HCT 116 出现频率最高（98 次，25.6\%），其次为肝癌细胞 HepG2（41 次，10.7\%）、黑色素瘤细胞 A375（23 次，6.0\%）和乳腺癌细胞 MCF-7（23 次，6.0\%），前 10 种细胞系合计覆盖 71.0\% 的样本。
组织来源以结肠（28.7\%）、肝脏（16.7\%）和脑组织（8.9\%）为主。

表面修饰类型共 43 种，其中聚糖修饰（PTR）占比最高（127 次，33.2\%），其后依次为聚糖蛋白复合物 PR（18 次，4.7\%）和 PAMAM+RGD+siRNA 复合修饰（15 次，3.9\%）。
合成方法以化学还原法为主（352 次，占 91.9\%），毒性检测方法中最常见的是 MTT 法（206 次，53.8\%）。

\begin{figure}[H]
  \centering
  \begin{minipage}{0.49\linewidth}
    \centering
    \begin{tikzpicture}
      \begin{axis}[
        acs xbar,
        width=\linewidth,
        height=5.8cm,
        xmax=215,
        xlabel={样本量},
        symbolic y coords={其他细胞系,MCF-7,A375,HepG2,HCT 116},
        ytick=data,
        y tick label style={font=\scriptsize},
        title={\textbf{(A)} 主要细胞系}
      ]
        \addplot+[fill=acsblue, draw=none] coordinates {(198,其他细胞系) (23,MCF-7) (23,A375) (41,HepG2) (98,HCT 116)};
      \end{axis}
    \end{tikzpicture}
  \end{minipage}
  \hfill
  \begin{minipage}{0.49\linewidth}
    \centering
    \begin{tikzpicture}
      \begin{axis}[
        acs axis,
        width=\linewidth,
        height=5.8cm,
        ymin=0,
        ymax=65,
        ylabel={样本占比（\%）},
        symbolic x coords={PTR,PR,PAMAM+RGD+siRNA,其他修饰},
        xtick=data,
        x tick label style={rotate=35, anchor=east, font=\tiny},
        nodes near coords,
        nodes near coords style={font=\tiny, yshift=1pt},
        title={\textbf{(B)} 表面修饰类型}
      ]
        \addplot+[only marks, acsorange, mark=square*, mark size=2pt] coordinates {(PTR,33.2) (PR,4.7) (PAMAM+RGD+siRNA,3.9) (其他修饰,58.2)};
      \end{axis}
    \end{tikzpicture}
  \end{minipage}
  \caption{主要细胞系与表面修饰类型的频率分布。}
  \label{fig:senp-categorical-frequency}
\end{figure}

因此，本研究的可视化分析并不只用于展示变量分布，而是与特征筛选和模型评价相结合：先通过统计图和表格检查样本覆盖范围、异常高剂量和类别集中度，再根据模型表现选择需要解释的关键特征，并在具备完整模型预测结果后利用 PDP、ICE 和双特征等高线图分析浓度、粒径、Zeta 电位、暴露时间、检测方法及材料化学增强特征等变量对细胞毒性预测的单独效应和交互效应。

